% Frontmatter — Top-tier rhythm. Title aligned with body claim
% (cover-bidder population, not cartel ringleaders). Abstract: 8 sentences.

\begin{frontmatter}

\title{Frequent Losers as a Cartel Screen: Detecting Cover Bidding Without Bid Microdata\texorpdfstring{\footnote{This version: May 2026. We thank seminar participants at INSPER for comments on earlier drafts; all remaining errors are our own. \textbf{Data and code:} the replication archive (R, Python, and \LaTeX{} sources) is available upon publication; CADE adjudication records and BEC contract awards are public. \textbf{Funding:} none. \textbf{Declarations of interest:} none.}}{}}

\author[insper]{Darcio Genicolo-Martins\corref{cor1}}
\ead{darciogm1@insper.edu.br}
\author[insper]{Paulo Furquim de Azevedo}
\affiliation[insper]{organization={INSPER},
  city={S\~ao Paulo}, country={Brazil}}
\cortext[cor1]{Corresponding author.}

\begin{abstract}
Most tools for detecting bid-rigging cartels need bid microdata,
which few jurisdictions audit at scale. A peculiar pattern shows
up in São Paulo's electronic procurement platform (BEC,
$\valYearStart$--$\valYearEnd$): out of $\valAlwaysLosers$ firms
that bid without ever winning, $\valFL$ participated in
$\valThresholdStat$ or more tender-items. We call this group
\emph{frequent losers}. Their behavior is hard to square with
expected-profit competition and reads naturally as the
cover-bidder side of bid-rigging arrangements. Comparing items
within the same product code, year, and procuring unit
($\valSampleN$ tender-items in total), the presence of a
frequent loser is associated with a $\valHeadlineRange$ higher
negotiated unit price across five estimators (full-sample OLS,
PBU fixed effects, IPW, CEM, cross-fit). The price gap shrinks
as buyers grow: from $\valOversightQonePct$ in the smallest
procuring-unit-size quartile to $\valOversightQfourPct$ in the
largest. The resulting $\valPBUgradientRatio\!\!\times$ ratio is
the strongest source of heterogeneity in the data, larger than
the preg\~ao--convite gap or any cross-sector spread. We read the
pattern through a separating-equilibrium framework with cover
bidders and genuine entrants, in which
$\log(1+\text{tenders\_count})$ is a sufficient ranking statistic
for the cover-bidder type and $\partial m^*/\partial \theta_k < 0$
matches the oversight gradient. Identification is descriptive by
design. Because the construct is computed entirely from
contract-award records, we translate it into a three-stage
administrative pathway---Screen, Triage, Forensic---that
oversight bodies can deploy where bid microdata are not centrally
collected.
\end{abstract}

\end{frontmatter}

\noindent\textbf{Keywords:} frequent losers, cartel screen, cover bidding, public procurement, bid rigging, separating equilibrium

\noindent\textbf{JEL No.:} D44, D73, H57, K21, L41

\bigskip
